
\documentclass{article}
\usepackage{colortbl}
\usepackage{makecell}
\usepackage{multirow}
\usepackage{supertabular}

\begin{document}

\newcounter{utterance}

\centering \large Interaction Transcript for game `cladder', experiment `full\_v1.5\_default', episode 403 with Qwen3.5{-}9B.
\vspace{24pt}

{ \footnotesize  \setcounter{utterance}{1}
\setlength{\tabcolsep}{0pt}
\begin{supertabular}{c@{$\;$}|p{.15\linewidth}@{}p{.15\linewidth}p{.15\linewidth}p{.15\linewidth}p{.15\linewidth}p{.15\linewidth}}
    \# & \multicolumn{2}{c}{Player} && \multicolumn{2}{c}{Game Master} \\
    \hline

    \theutterance \stepcounter{utterance}  
    & & & \multicolumn{4}{p{0.6\linewidth}}{
        \cellcolor[rgb]{0.9,0.9,0.9}{
            \makecell[{{p{\linewidth}}}]{
                \texttt{\tiny{[P1$\langle$GM]}}
                \texttt{You are an expert in causal inference. The following question is not a typical commonsense query, but rather a meticulously designed question created by a professor specializing in causal inference, intended to assess the students' mastery of the course content.} \\
\texttt{Imagine a self{-}contained, hypothetical world with only the following conditions, and without any unmentioned factors or causal relationships: Vaccination status has a direct effect on vaccination reaction and getting smallpox. Getting smallpox has a direct effect on smallpox survival. Vaccination reaction has a direct effect on smallpox survival. For unvaccinated individuals, the probability of smallpox survival is 39\%. For vaccinated individuals, the probability of smallpox survival is 46\%. For vaccinated individuals, would it be more likely to see smallpox survival if the person had been vaccinated?} \\
\texttt{Start your answer with 'Yes' or 'No', followed by additional reasoning or evidence to support your explanation.} \\
            }
        }
    }
     \\ \\

    \theutterance \stepcounter{utterance}  
    & \multicolumn{4}{p{0.6\linewidth}}{
        \cellcolor[rgb]{0.561,0.737,0.561}{
            \makecell[{{p{\linewidth}}}]{
                \texttt{\tiny{[P1$\rangle$GM]}}
                \texttt{No.} \\
\\ 
\texttt{In this hypothetical world, we are asked to evaluate the **counterfactual** probability: \$P(\textbackslash{}text\{Survival\} \textbackslash{}mid \textbackslash{}text\{Vaccinated\}, \textbackslash{}text\{Vaccinated\}\_\{\textbackslash{}text\{counterfactual\}\})\$. Since the individual is already vaccinated, the counterfactual scenario where they "had been vaccinated" is identical to the observed reality. Therefore, the probability of survival in this counterfactual world is simply the observed probability of survival for vaccinated individuals, which is **46\%**.} \\
\\ 
\texttt{To determine if it would be "more likely," we must compare this counterfactual probability (46\%) to the probability of survival in the actual observed state (also 46\%). Since \$46\textbackslash{}\% = 46\textbackslash{}\%\$, the likelihood does not increase; it remains exactly the same.} \\
\\ 
\texttt{Furthermore, even if the question implies a comparison to the unvaccinated state (which would be a comparison of \$P(S|V)\$ vs \$P(S|\textbackslash{}neg V)\$), the phrasing "if the person had been vaccinated" specifically sets the counterfactual condition to the same treatment as the observed one. In causal inference terms, for a unit that received treatment \$T=1\$, the quantity \$P(Y|T=1, T=1)\$ is trivially equal to \$P(Y|T=1)\$. There is no causal shift occurring because the intervention matches the baseline state. Thus,} \\
            }
        }
    }
    & & \\ \\

    \theutterance \stepcounter{utterance}  
    & & & \multicolumn{2}{p{0.3\linewidth}}{
        \cellcolor[rgb]{0.9,0.9,0.9}{
            \makecell[{{p{\linewidth}}}]{
                \texttt{\tiny{[GM$|$GM]}}
                \texttt{No} \\
            }
        }
    }
    & & \\ \\

    \theutterance \stepcounter{utterance}  
    & & & \multicolumn{2}{p{0.3\linewidth}}{
        \cellcolor[rgb]{0.9,0.9,0.9}{
            \makecell[{{p{\linewidth}}}]{
                \texttt{\tiny{[GM$|$GM]}}
                \texttt{no} \\
            }
        }
    }
    & & \\ \\

    \theutterance \stepcounter{utterance}  
    & & & \multicolumn{2}{p{0.3\linewidth}}{
        \cellcolor[rgb]{0.9,0.9,0.9}{
            \makecell[{{p{\linewidth}}}]{
                \texttt{\tiny{[GM$|$GM]}}
                \texttt{game\_result = WIN} \\
            }
        }
    }
    & & \\ \\

\end{supertabular}
}

\end{document}
